% ---------------------------------------------------------------------
% II. BACKGROUND AND RELATED WORK (about 0.5 page)
% Owner: Rory Linnane
%
% ---------------------------------------------------------------------
\section{Background and Related Work}\label{sec:background}

\subsection{Bias in Pretraining Data Filters}
Large language models are trained on web-scale corpora filtered through pipelines that silently encode value judgements.
Dodge et al.\ showed that C4's word blocklist disproportionately removes text by and about LGBTQ+ people and in African American English (AAE)~\cite{dodge2021documenting}.
Quality filters compound this: Gururangan et al.\ found that a replica of GPT-3's quality filter assigns higher scores to text from wealthier, urban and more educated sources~\cite{gururangan2022whose}.
Toxicity classifiers exhibit the same pattern: Sap et al.\ showed hate-speech detectors flag AAE at significantly higher rates than equivalent white-aligned English~\cite{sap2019risk}, and Mendu et al.\ found that a transformer-based filter misclassifies news articles that \emph{report} on harmful content as harmful themselves~\cite{mendu2025safer}.
Fairness-oriented filters are not immune: Udagawa et al.\ found a regard classifier removes sentences that report real-world events involving marginalised groups~\cite{udagawa2025bias}.
Identity-term templates offer a controlled way to probe these failures.
Dixon et al.\ introduced the minimal-pair method, which substitutes one identity word for a neutral control (e.g.\ ``my neighbour is \emph{gay}'' vs.\ ``my neighbour is \emph{tall}'') to isolate unintended classifier bias~\cite{dixon2018measuring}.
Borkan et al.\ extended this work to real comments, releasing the identity-labelled Civil Comments data on which Detoxify \texttt{unbiased} was later trained~\cite{borkan2019nuanced}.
However, across all filter types, these audits measure only aggregate removal rates by group, and none explains why an individual sentence was removed.

\subsection{Feature Attribution for Text}
Post-hoc feature attribution assigns each input token a score reflecting its contribution to the model's prediction, without modifying the model itself.
SHAP grounds attribution in cooperative game theory: Lundberg and Lee showed that the Shapley value, the unique attribution satisfying efficiency, symmetry and the null-player axiom, can be computed for any model via a unified additive framework~\cite{lundberg2017unified}.
For text, the SHAP library uses a partition approximation; exact Shapley enumeration is feasible only for short sequences, which allows it to serve as a reference (Sec.~\ref{sec:method}).
LIME takes a complementary approach: Ribeiro et al.\ fit a locally weighted sparse linear surrogate to random perturbations of the input, measuring local fidelity through the surrogate's $R^2$~\cite{ribeiro2016why}.
Both methods are post-hoc, local (one prediction at a time) and model-agnostic: they treat the classifier as a black box and require only query access~\cite{munn2022xai}.

\subsection{Evaluating Explanations: Faithfulness vs.\ Agreement}
Jacovi and Goldberg define an explanation as \emph{faithful} if it correctly describes the model's actual reasoning process, and distinguish this from whether it looks plausible to a human~\cite{jacovi2020faithfully}.
Deletion-based metrics operationalise faithfulness: DeYoung et al.'s ERASER benchmark defines \emph{comprehensiveness} (the score drops when the top-$k$ attributed tokens are deleted) and \emph{sufficiency} (the score is preserved when only those tokens are kept)~\cite{deyoung2020eraser}.
These metrics have a known limitation: deleting words produces inputs outside the model's training distribution, so score changes may partly reflect distributional shift rather than token importance~\cite{hase2021ood}.
More broadly, post-hoc explainers are imperfect.
Krishna et al.\ showed that SHAP and LIME frequently disagree on feature rankings for the same input, so agreement between them does not imply either is correct~\cite{krishna2022disagreement}.
Slack et al.\ demonstrated that explanations can be adversarially manipulated to conceal discriminatory model behaviour~\cite{slack2020fooling}.
Ngueajio et al.\ found that the highest-accuracy classifier is not necessarily the one whose explanations are most faithful~\cite{ngueajio2025decoding}.
These evaluations were conducted on general classification tasks; to our knowledge, no prior work has evaluated attribution faithfulness in pretraining filter audits against references whose correct answers are known independently.
